% year: תש"פ
% type: בוחן
% semester: א
% points: 20
% source: exams/בחנים/תש״פ, בוחן, 9141.pdf
% number: 5
% categories: continuity, derivatives, multiple-choice

\begin{question}
הפונקציה $f(x)$ רציפה בנקודה $x=x_0$. איזו טענה נכונה ?
\begin{enumerate}
  \item $f(x)$ גזירה בנקודה $x=x_0$
  \item קיים גבול $\displaystyle\lim_{x\to x_0}\frac{f(x)-f(x_0)}{x-x_0}$
  \item קיים גבול $\displaystyle\lim_{x\to x_0}\left(f(x)-f(x_0)\right)$
  \item קיים גבול $\displaystyle\lim_{h\to 0}\frac{f(x_0+h)-f(x_0)}{h}$
  \item $f'(x)$ רציפה בנקודה $x=x_0$
\end{enumerate}
\end{question}

\begin{hint}
רציפות אינה גוררת גזירות. חשבו על $f(x)=|x|$ בנקודה $x_0=0$.
\end{hint}

\begin{solution}
\textbf{התשובה הנכונה: (ג)}.

לפי ההגדרה, $f$ רציפה ב-$x_0$ אם ורק אם $\lim_{x\to x_0}f(x)=f(x_0)$. לפי אריתמטיקה של גבולות (גבול של קבוע הוא הקבוע עצמו) זה שקול לכך ש-
\[
\lim_{x\to x_0}\bigl(f(x)-f(x_0)\bigr)=f(x_0)-f(x_0)=0 ,
\]
ובפרט הגבול בטענה (ג) קיים (ושווה ל-0).

\textbf{למה האחרות שגויות:} נתבונן ב-$f(x)=|x|$ וב-$x_0=0$. הפונקציה רציפה ב-0, אבל
\[
\frac{f(0+h)-f(0)}{h}=\frac{|h|}{h}=\begin{cases}1,& h>0\\ -1,& h<0\end{cases}
\]
ולכן הגבולות החד-צדדיים של מנת ההפרשים הם $1$ ו-$-1$ והגבול אינו קיים. מכאן:
\begin{itemize}
  \item (ב) ו-(ד) שגויות: אלה שתי צורות כתיבה של הגדרת הנגזרת $f'(x_0)$, והראינו שהגבול לא קיים.
  \item (א) שגויה: גזירות ב-$x_0$ פירושה בדיוק קיום הגבול ב-(ב)/(ד).
  \item (ה) שגויה: $f'(0)$ כלל אינה מוגדרת, ולכן $f'$ אינה רציפה ב-0. (יתרה מזו, גם פונקציה גזירה יכולה להיות בעלת נגזרת לא רציפה, למשל $x^2\sin\frac1x$ עם $f(0)=0$.)
\end{itemize}
\end{solution}
